\section{$3$-dynamic coloring of $C_p(1,3)$}

\begin{theorem}\label{thm:circulant-r3}
Let $p\ge12$. Then
\[
\chi_3^d(C_p(1,3))=
\begin{cases}
4, & p\equiv0\pmod{12},\\
5, & \text{otherwise}.
\end{cases}
\]
\end{theorem}

\begin{proof}
Suppose first that $12\mid p$. The cyclic color word
$
(010232101323)^{p/12}
$
defines a proper $3$-dynamic $4$-coloring. Hence
$\chi_3^d(C_p(1,3))=4$ by
Equation~\ref{eq:absolute-bounds}.

We next give a $3$-dynamic $5$-coloring for every $p\ge12$.
Let $A=01234$ and write $p=5q+r$, where $0\le r\le4$.
Consider the following cyclic color words:
\[
\begin{array}{c|ccccc}
r&0&1&2&3&4\\ \hline
\phi&
A^q&
A^{q-1}102431&
A^q01&
A^{q-1}01314203&
A^q1203
\end{array}
\]
A direct verification shows that each word defines a proper
$3$-dynamic $5$-coloring. Thus
$
\chi_3^d(C_p(1,3))\le5
$
for every $p\ge12$.

It remains to show that a $3$-dynamic $4$-coloring can exist only if
$12\mid p$. Let $\phi$ be such a coloring. Since $\phi$ is proper,
the neighbors of every vertex use exactly the three colors different
from the color of the vertex. Thus exactly two of its four neighbors
receive the same color.

Define
\[
s_i=
\begin{cases}
1,&\phi(i-1)=\phi(i+1),\\
0,&\phi(i-1)\neq\phi(i+1).
\end{cases}
\]
We first observe that the repeated color in $N(i)$ can only occur on
one of the pairs
$
\{i-3,i-1\},\qquad
\{i-1,i+1\},\qquad
\{i+1,i+3\}.
$
Indeed, if $\phi(i-3)=\phi(i+1)$, then properness forces
$\phi(i+2)=\phi(i)$, and the $3$-dynamic condition at $i-1$ forces
one of $i-4,i-2$ to receive the color of $i-3$, contradicting
properness. The case $\phi(i-1)=\phi(i+3)$ is symmetric. If
$\phi(i-3)=\phi(i+3)$, then properness forces
$
\phi(i-2)=\phi(i+2)=\phi(i),
$
contradicting the $3$-dynamic condition at $i-1$.

Consequently, for every $i$, exactly one of
\[
s_{i-2},\ s_i,\ s_{i+2}
\]
is equal to $1$. Moreover, two consecutive $s_i$ cannot both be $1$.
Indeed, if $s_i=s_{i+1}=1$, a direct reconstruction using the
$3$-dynamic condition around $i$ and $i+1$ forces
$\phi(i+4)=\phi(i+5)$, contradicting properness.

Choose $i$ with $s_i=1$. Then by the previous
observations, $s_{i-6}=s_{i+6}=s_{i-3}=s_{i+3}=1$ and every other $s_k$ between $i-6$ and $i+6$ equals 0.
This implies $3\mid p$.
Reconstructing $\phi$ yields
$\phi(i-6)=\phi(i+6)$.
Repeating the same argument gives
\[
\phi(i-6)=\phi(i-6+12k) \neq \phi(i)
\]
for every $k$, and hence, as $p\equiv0\pmod {12}$.
\end{proof}